\documentclass[10pt,a4paper]{amsart}
\usepackage[margin=19mm]{geometry}
\usepackage{amsmath,amssymb,mathtools}
\usepackage[scr=rsfs,cal=euler]{mathalfa}
\usepackage{xfrac}
\usepackage[unicode,hidelinks,bookmarks=false]{hyperref}
\newcommand{\ie}{i.e.\ }
\newcommand{\cf}{cf.\ }
\newcommand{\resp}{respectively\ }
\newcommand{\Subsection}{\subsection}
\providecommand{\nobreakdash}{\nobreak\hbox{-}}
\setlength{\emergencystretch}{2em}
\multlinegap0pt
\usepackage{bm}
\usepackage{bbm}
\usepackage{dsfont}
\usepackage{xspace}
\usepackage{cancel}
\usepackage{mathtools}
\usepackage{amsthm}
\makeatletter
\@ifundefined{theorem}{\newtheorem{theorem}{Theorem}}{}
\makeatother
\makeatletter
\expandafter\ifx\csname proof\endcsname\relax
\newenvironment{proof}[1][Proof]
%{\ \rule{0.5em}{0.5em}}

\fi
\makeatother
\title[Dedekind-finite cardinals having countable partitions]{Dedekind-finite cardinals having countable partitions}
\author{Supakun Panasawatwong, J K Truss}
\date{Source: arXiv:2204.05700v1 (CC BY 4.0)}
\begin{document}
\maketitle

\noindent\emph{Source and licence.} Dedekind-finite cardinals having countable partitions. Version arXiv:2204.05700v1 (CC BY 4.0), distributed under the Creative Commons Attribution 4.0 International licence, as recorded in the arXiv licence metadata for this exact version and shown on the abstract page https://arxiv.org/abs/2204.05700v1. Full rights notice: licenses/arxiv-2204.05700-CC-BY-4.0.txt.\par\smallskip

\noindent\textit{Editorial scope.} Counted unit: the Theorem whose source label is \texttt{3.2} in the article (source0.tex lines 377--379, source label \texttt{3.2}), delivered together with the article's own complete original proof of it (source0.tex lines 381--466, one \texttt{proof} environment of 86 lines). No mathematics is written, repaired or re-derived: statement and proof are the article's own text line for line. Required context retained uncounted: none. Every definition, display and label of those lines is retained unchanged. Omitted from this package and not redistributed: 1--376, 380--380, 467--801. Nothing inside a retained excerpt is omitted or altered: every retained body is delivered exactly as the article's lines are written, so no omission receipt of any kind accompanies this package. Citations: the retained excerpts carry no citation command at all, so no bibliography entry of the article is transcribed and no reference is redistributed. Reference adaptation: none was needed and none was made; every reference inside the delivered unit resolves inside it, so no unresolved reference or citation is produced. Printed slips kept literal: no printed word, symbol, hypothesis or number of the counted statement or of its proof was altered. Layout and class: the article's own class is \texttt{amsart} with packages including amssymb, amsfonts; the wrapper is the lane's pinned \texttt{amsart} editorial wrapper at 10pt on A4, which restates the theorem-like environments the retained text needs and adds the article's own macro definitions that the retained text needs (none), with extra wrapper packages (bm, bbm, dsfont, xspace, cancel, mathtools). The delivered numbering is the unit's own. Assets: no retained span contains a figure environment, a graphics-inclusion command, a PDF-page inclusion, a table or a TikZ picture; no image file is copied into this package, the package declares no asset dependency and no image file is redistributed with it. Licence: version arXiv:2204.05700v1 of this article is distributed under the Creative Commons Attribution 4.0 International licence, as recorded in the arXiv licence metadata for this exact version and shown on the abstract page https://arxiv.org/abs/2204.05700v1. A full rights notice is delivered at licenses/arxiv-2204.05700-CC-BY-4.0.txt. Characters: every retained excerpt is pure ASCII, so the delivered engine, XeLaTeX with its default Latin Modern text font, reports no missing character.
\begin{theorem}  \label{3.2} For any countable weakly $2$-transitive trees $T_1$, $T_2$ in which all maximal chains
of $T_1^+$, $T_2^+$ are dense, having distinct infinite codes, $|U_{T_1}|$ and $|U_{T_2}|$ are $\equiv$-inequivalent 
members of $\Delta_5 \setminus \Delta_4$.   \end{theorem}

\begin{proof}  The point of the restriction we have made on our weakly 2-transitive trees is that for any
ramification point $r$ and $x, y > r$ in $T$, there is an automorphism of $T$ fixing $r$ and taking $x$ to $y$.

Suppose for a contradiction that $|U_{T_1}| \equiv |U_{T_2}|$. In ${\mathfrak N}_{T_1}$ there must therefore be a 
structure $\prec$ on $U_{T_1}$ such that $(U_{T_1}, \prec)$ in ${\mathfrak N}_{T_1}$ is elementarily equivalent to 
$(U_{T_2}, <)$ in ${\mathfrak N}_{T_2}$. Let $A$ be a finite support for $(U_{T_1}, \prec)$ in ${\mathfrak N}_{T_1}$. 
Thus any automorphism of $(U_{T_1}, <)$ that fixes $A$ pointwise also fixes $\prec$. We may assume that $A$ is a 
subtree of $U_{T_1}$.

In order to handle this situation, we use the following notation, where $A$ is a finite subset of $T$ and $[A]$ 
is its definable closure, in this case just the lower subsemilattice generated by $A$. If $a$ is maximal in $A$ 
(or if $A = \emptyset$, $a$ can be $-\infty$), we let $U_a = \{x \in T: a < x\}$, and if $a < b$ are consecutive 
members of $[A]$, then $L_{ab} = \{x \in T: a < x < b\}$ (the linear piece between $a$ and $b$; again here $a$ 
can be $-\infty$ in the case when $b$ is minimal in $A$), and $S_{ab} = \{x \in T: a < x, b \not \le x\}$ (the 
corresponding side piece). Now $L_{ab}$ is linear, and we may also consider $L_{ab}^+$, which is the same thing, 
but calculated in $T^+$, that is, including ramification points too, which are coloured according to which orbit 
they lie in. If $x \in L_{ab}^+$, then $C_x = \{y \in T: b \wedge y = x\}$, which is the set of points which 
branch off from $L_{ab}^+$ at $x$. This is a union of cones at $x$.

A key remark is that the pointwise stabilizer $G_A$ acts transitively on each $U_a$, $L_{ab}$, and $C_x$ (here heavily 
using the assumption that the maximal chains in $T^+$ are dense). But more is true. 
In fact for any $b, c \in U_a$, there is an automorphism of $T$ taking $b$ to $c$ which fixes $A$ pointwise, {\em and} has
support contained in $U_a$, with a similar statement holding for $C_x$. The analogous statement also holds for
weak 2-transitivity. For instance, if $a < b < c$ and $a < d < e$, then by weak 2-transitivity there is an 
automorphism $f$ fixing $a$ and taking $b$ to $d$, and applying weak 2-transitivity again, there is an automorphism
$g$ fixing $d$ and taking $f(c)$ to $e$, and $f$, $g$ may be chosen fixing all points below $a$ and $d$ respectively.
Then $gf$ fixes $a$ and takes $b$ to $d$ and $c$ to $e$.

Although not actually needed, we can describe the orbits of $G_A$. They are of the form $U_a$, and $L_{ab}$, as well 
as the union of the cones in $S_{ab}$ which meet $L_{ab}$ at points lying in some orbit of $G_A$ on $L_{ab}^+$.

Now being a tree with dense chains is first order expressible, and so this must be true of $(U_{T_1}, \prec)$. We shall show
that on each set of the form $U_a$ or $C_x$, $\prec$ is equal to $<$ or the empty set (in which case it would be an antichain, 
i.e. pairwise $\prec$-incomparable). We just treat $U_a$, as $C_x$ is similar.

First consider a maximal point $a$ of $A$ (and let $a$ have a default value of $-\infty$ if $A = \emptyset$). First we note
that if $b, c \in U_a$ are $<$-incomparable, then they are also $\prec$-incomparable. This is because there is an 
element of $G_A$ swapping $b$ and $c$. Since $G_A$ preserves $\prec$, $b \prec c \Leftrightarrow c \prec b$, and since
$\prec$ is antisymmetric, $\neg b \prec c \wedge \neg c \prec b$. 

Next we show that there cannot be points $b, c \in U_a$ such that $b < c$ and $c \prec b$. For, if so, by weak
2-transitivity of $G_A$ on $U_a$, for $x, y \in U_a$, $x < y \Rightarrow y \prec x$. Choose $z > a$ and incomparable 
$x, y > z$. Thus $x, y \prec z$. By the previous paragraph, $x$ and $y$ are also $\prec$-incomparable, so this 
contradicts $(U_{T_1}, \prec)$ a tree. 

Next suppose that there are $b, c \in U_a$ such that $b < c$, and $b$ and $c$ are $\prec$-incomparable. By 
weak 2-transitivity of $U_{T_1}$, $a < x < y$ implies that $x$, $y$ are $\prec$-incomparable. But also as shown above, if
$b$ and $c$ are $<$-incomparable, they are also $\prec$-incomparable, and we deduce that $\prec$ is the empty relation
on $U_a$.

Otherwise, it follows that for $x, y \in U_a$, $x < y \Rightarrow x \prec y$. Since also $x, y$ $<$-incomparable
$\Rightarrow x, y \prec$-incomparable and $y < x \Rightarrow y \prec x$, we deduce that $x < y \Leftrightarrow x \prec y$. 
In other words, $<$ and $\prec$ agree on $U_a$.

Now let us consider $S_{ab}$ where $a < b$ are consecutive members of $A$. First note that by a proof similar to the above,
if $x, y$ are $<$-incomparable members of $C_c$ for $a < c < b$, then they are also $\prec$-incomparable, since there is a 
member of $G_A$ which swaps $x$ and $y$ (and fixes $c$). We can deduce that if $a < c < d < b$, and $x \in C_c$, $y \in C_d$,
then $x$ and $y$ are $\prec$-incomparable. For suppose that $x \prec y$ (a similar argument applying if $y \prec x$). Choose
$x' \in C_c$ which is $<$-incomparable with $x$. Then by what we have just shown, $x$, $x'$ are $\prec$-incomparable. 
Furthermore, there is a member of $G_A$ which swaps $x$ and $x'$, and having support contained in $C_c$, and so this also
fixes $y$. Since $G_A$ preserves $\prec$, $x' \prec y$, but this contradicts $\prec$ a tree.

It follows that any $\prec$-comparabilities in $S_{ab}$ must hold between members of the same $C_c$. Suppose then that
$x < y$ and $x \prec y$ hold for some $x, y \in C_c$. As $G_A$ acts weakly 2-transitively on $C_c$, $x < y \Rightarrow x \prec y$. 
Since we have also shown that $x$, $y$ $<$-incomparable implies that they are $\prec$-incomparable, we deduce that $<$ and $\prec$
agree on $C_c$. 

The main remaining point is to show that there is some $a$ or $c$ such that $\prec$ is non-empty on $U_a$ or $C_c$. Suppose not. Then
each $U_a$ is an antichain, and since the members of distinct $C_c$s for $a < c < b$ are incomparable, each $S_{ab}$ is also an
antichain. Consider the possible relationship between $a \in A$ and members of $U_b$ under $\prec$. The group $G_A$ acts transitively 
on $U_b$, and so if any member of $U_b$ is $\prec$-less than $a$, they all are, which violates $\prec$ a tree. Similarly, for members
of $C_c$. The same argument applies to the possible relationship between members of distinct $U_a$s, or $C_c$s, or between members 
of $U_a$ and $C_c$. We deduce that the union of all the $U_a$s and $S_{ab}$s is an antichain, and all its members are either 
incomparable with every member of $A$, or above some member of $a$. However, since $(U_{T_1}, \prec)$ is a tree satisfying the
same first order sentences as $(U_{T_1}, <)$, all its chains are infinite, so this is impossible.

We deduce that there is $a$ or $c$ such that $<$ and $\prec$ agree on $U_a$ or $C_c$. We just treat the first case, and show that 
for any $b \in U_a$, for $x \in T_1$, $b < x \Leftrightarrow b \prec x$. From left to right is already known. Suppose for a 
contradiction that there are some $b \in U_a$ and $x$ such that $b \prec x$ but not $b < x$. Since $<$ and $\prec$ agree on $U_a$, 
$x \not \in U_a$. Choose $c \in U_a$ incomparable with $b$. Then there is a member of $G_A$ taking $b$ to $c$ and with support 
contained in $U_a$. This preserves $\prec$, and fixes $x$, and hence $c \prec x$. Since $\prec$ is a tree, $b$ and $c$ are 
$\prec$-comparable, and as they both lie in $U_a$, they are also $<$-comparable, which gives a contradiction.

Finally we choose any $b \in U_a$, and observe that the ramification orders arising above $b$ in $<$ and $\prec$ are equal to those
arising in $T_1$ and $T_2$ respectively. Since these are determined by first order formulae, it follows that $|U_{T_1}|$ and $|U_{T_2}|$ 
have equal codes after all.  \end{proof}

\end{document}
